\chapter{Untersuchungsmethode}
\label{chap:methode}

\section{Vorgehensweise und Methode}

In der vorliegenden Arbeit wird ein theoretischer Ansatz verfolgt: Es werden keine eigenen Primärdaten erhoben, sondern wissenschaftliche Fachliteratur, amtliche Dokumente und institutionelle Berichte ausgewertet und verglichen. Diese Methode wurde gewählt, da sie es ermöglicht, die Forschungsfrage auf Basis veröffentlichter Studien und offizieller Dokumente für mehrere Regionen zu beantworten. Eine eigene Erhebung wäre hierfür nicht geeignet, da Umfragen weder die geografische Breite mehrerer Regionen abdecken noch verlässliche Aussagen über die wirtschaftspolitischen Entwicklungen ganzer Regionen liefern können. Als Auswahlkriterien galten thematische Relevanz, Aktualität (2018–2026) sowie die Herkunft aus wissenschaftlichen Datenbanken oder offiziellen Institutionen. Für jede der drei Regionen wurden relevante Studien und Dokumente ausgewertet und mithilfe eines komparativen Ansatzes gegenübergestellt, um Gemeinsamkeiten und Unterschiede herauszuarbeiten~\cite{Yin2018}.

\section{Fallauswahl}

Für den Vergleich wurden El Salvador, die Vereinigten Arabischen Emirate (VAE) und die Europäische Union (EU) ausgewählt, da sie drei grundlegend verschiedene staatliche Strategien repräsentieren. So kann untersucht werden, ob trotz unterschiedlicher Ausgangsbedingungen gemeinsame Hindernisse erkennbar sind~\cite{Flick2017}:

\begin{itemize}
    \item \textbf{El Salvador} -- vollständige staatliche Akzeptanz als gesetzliches
    Zahlungsmittel~\cite{Alvarez2023}
    \item \textbf{VAE} -- wirtschaftsorientierte Integration durch regulierten
    Markt~\cite{VARA2025}
    \item \textbf{EU} -- schrittweise Regulierung zum Schutz der
    Verbraucher~\cite{MiCA2023}
\end{itemize}
